% The next line makes it so it compiles the main file if I hit the compile button
% when editing this file in TeXworks.
% !TEX root = ../ActiveIntroToDiscreteMathAndAlgorithms.tex
\chapter{Counting}\label{ch:count}
In this chapter we provide a very brief introduction to a field called {\em combinatorics}.
\index{combinatorics}
We are actually only going to scratch the surface of this very broad and deep subfield
of mathematics and theoretical computer science.  We will focus on a subfield of
combinatorics that is sometimes called {\em enumeration}.  That is, we will mostly 
concern ourselves with how to count things.

It turns out that combinatorial problems are notoriously deceptive.
Sometimes they can seem much harder than they are, and at other times they seem easier than
they are.  In fact, there are many cases in which one combinatorial problem will be relatively
easy to solve, but a very closely related problem that seems almost identical will be 
very difficult to solve.

When solving combinatorial problems, you need to make sure you fully understand what is 
being asked and make sure you are taking everything into account appropriately.
I used to tell students that combinatorics was easy.  I don't say that anymore.  In some
sense it is easy.  But it is also easy to make mistakes.

\section{The Sum and Product Rules}\label{sec:multSum}
We begin our study of combinatorial methods with the following two
fundamental principles.  They are both pretty intuitive.  The only difficulty
is realizing which one applies to a given situation.  If you have a good
understanding of what you are counting, the choice is generally pretty clear.

\begin{thm}[Sum Rule] \index{sum rule}
\label{rul:sum}Let $E_1, E_2, \ldots, E_k$, be pairwise finite disjoint sets. Then
$$|E_1 \cup E_2 \cup \cdots \cup E_k| = |E_1| + |E_2| + \cdots + |E_k|.$$

Another way of putting the sum rule is this:  
If you have to accomplish some task and you can do it in one of $n_1$ ways, or
one of $n_2$ ways, etc., up to one of $n_k$ ways, and none of the ways of doing the
task on any of the list are the same, then there are
$n_1+n_2+\cdots+n_k$ ways of doing the task.
\end{thm}


\begin{exa}
I have 5 brown shirts, 4 green shirts, 10 red shirts, and 3 blue shirts.  How many choices
do I have if I intend to wear one shirt?
\begin{sol}
Since each list of shirts is independent of the others, I can use the sum rule.
Therefore I can choose any of my $5+4+10+3=22$ shirts.
\end{sol}
\end{exa}

\begin{exa}
How many ordered pairs of integers $(x, y)$ are there such that $0<
|xy| \leq 5$?
\begin{sol} Let $E_k = \{(x, y) \in \BBZ^2: |xy| = k\}$ for $k = 1,
\ldots , 5$. Then the desired number is 
$$|E_1| + |E_2| + \cdots + |E_{5}|.$$ 
We can compute each of these as follows:
$$\begin{array}{lll}E_1 & = & \{(-1,-1), (-1, 1), (1, -1), (1,1)\} \\
E_2 & = & \{(-2,-1), (-2, 1), (-1,-2), (-1, 2), (1, -2), (1,2), (2,-1), (2,1)\} \\
E_3 & = & \{(-3,-1), (-3, 1), (-1,-3), (-1, 3), (1, -3), (1,3), (3,-1), (3,1)\} \\
E_4 & = & \{(-4,-1), (-4, 1), (-2, -2), (-2, 2),  (-1,-4), (-1, 4), (1, -4), \\
&&\, \, (1,4), (2,-2),(2,2),  (4,-1), (4,1)\} \\
E_5 & = & \{(-5,-1), (-5, 1), (-1,-5), (-1, 5), (1, -5), (1,5), (5,-1), (5,1)\} \\
  \end{array}$$
  The desired number is therefore $4+ 8 + 8 + 12 + 8 = 40$.
\end{sol}
\end{exa}

\begin{exer}
For dessert you can have cake, ice cream or fruit.  There are $3$ kinds of cake,
$8$ kinds of ice cream and $5$ different of fruits.  How many choices do you have
for dessert?

\noindent
\answerline
\begin{answer}
Nobody in their right mind will choose fruit if cake and ice cream are
available, so there are $3+8=11$ choices.
Just kidding.  There are really $3+8+5=16$ different choices.
\end{answer}
\end{exer}

\begin{thm}[Product Rule]\index{product rule}
\label{rul:product} Let  $E_1, E_2, \ldots, E_k$, be finite  sets. Then
$$|E_1 \times E_2 \times \cdots \times E_k| = |E_1| \cdot |E_2|  \cdots  |E_k|.$$

Another way of putting the product rule is this:
If you need to accomplish some task that takes $k$ steps, and there are $n_1$ ways of
accomplishing the first step, $n_2$ ways of accomplishing the second step, etc., 
and $n_k$ ways of accomplishing the $k$th step, then there are $n_1n_2\cdots n_k$
ways of accomplishing the task.
\end{thm}


\begin{exa}
I have 5 pairs of socks, 10 pairs of shorts, and 8 t-shirts.  How many choices do I
have if I intend to wear one of each?
\begin{sol}
I can think of choosing what to wear as a task broken into 3 steps:
I have to choose a pair of socks (5 ways), a pair of shorts (10 ways), and finally a t-shirt 
(8 ways).  Thus I have $5\times 10\times 8= 400$ choices.
\end{sol}
\end{exa}


\begin{exer}
If license plates are required to have 3 letters followed by 3 digits, 
how many license plates are possible?

\noindent
\answerline
\begin{answer}
There are 26 choices for each of the first three characters, 
and 10 choices for each of the final three characters.  
Therefore, there are $26^3\cdot10^3$ possible license plates.
\end{answer}
\end{exer}


\begin{exa}\label{exa:positive_divisors}
The positive divisors of $400$ are written in increasing order
$$1,2,4,5,8, \ldots, 200, 400.
$$
How many integers are there in this sequence? 
How many of the divisors of $400$ are perfect squares?
\begin{sol} Since $400 = 2^4 \cdot 5^2$, any positive divisor of $400$
has the form $2^a5^b$ where $0 \leq a \leq 4$ and $0\leq b \leq 2$.
Thus there are $5$ choices for $a$ and $3$ choices for $b$ for a
total of $5\cdot 3 = 15$ positive divisors.

To be a perfect square, a positive divisor of $400$ must be of the
form $2^\alpha 5^\beta$ with $\alpha \in \{0, 2, 4\}$ and $\beta \in
\{0, 2\}$. Thus there are $3\cdot 2 = 6$ divisors of $400$ which are
also perfect squares.
\end{sol}
\end{exa}

It is easy to generalize Example \ref{exa:positive_divisors} 
to obtain the following theorem.

\begin{thm}\label{thm:positive_divisors}
Let the positive integer $n$ have the prime factorization 
$$n = p_1^{a_1}p_2^{a_2}\cdots p_k^{a_k},
$$where the $p_i$ are distinct primes, and the $a_i$ are integers
$\geq 1$. If $d(n)$ denotes the number of positive divisors of $n$,
then $$d(n) = (a_1 + 1)(a_2 + 1)\cdots (a_k + 1).   $$
\end{thm}
\begin{exer}
Prove Theorem~\ref{thm:positive_divisors}.
(Hint: Follow the idea from Example~\ref{exa:positive_divisors}.)
\vspace*{3.5in}
\begin{answer}
Every divisor of $n$ is of the form $p_1^{b_1}p_2^{b_2}\cdots p_k^{b_k}$,
where $0\leq b_1\leq a_1$, $0\leq b_2\leq a_2$, $\ldots$, $0\leq b_k\leq a_k$.
(We could also write this as $0\leq b_i\leq a_i$ for $0\leq i\leq k$.)
Therefore there are $a_1+1$ choices for $b_1$, 
$a_2+1$ choices for $b_2$, all the way through
$a_k+1$ choices for $b_k$.  
Since each of the $b_i$s are independent of each other,
the product rule tells us that the number of divisors of $n$ is
$(a_1 + 1)(a_2 + 1)\cdots (a_k + 1)$.
\end{answer}
\end{exer}

\begin{ques}
Whether or not you realize it, you used the fact that the $p_i$ were distinct primes
in your proof of Theorem~\ref{thm:positive_divisors} (assuming you did the proof
correctly).  Explain where that fact was used (perhaps implicitly).

\noindent
\answerlinetwo
\begin{answer}
Unless the $p_i$ are distinct, the $b_i$s are not independent of each other.
In other words, if the $p_i$s are distinct, then each different choice of the $b_i$s
will produce a different number.  But this is not the case if the $p_i$s are not 
distinct.  For instance, if we write $32=2^32^2$, we can get the factor $4$ as
$2^22^0$, $2^12^1$, or $2^02^2$.  Clearly we would count $4$ three times and
would obtain the incorrect number of divisors.
\end{answer}
\end{ques}

\begin{exa}
What is the value of $sum$ after each of the following segments of code?

\noindent
\begin{minipage}[t]{.45\textwidth}
{\small 
\begin{verbatim}
    int sum=0;
    for(int i=0;i<n;i++) {
        for(int i=0;i<m;i++) {
            sum = sum + 1;
        }
    }
\end{verbatim}
}
\end{minipage}
\hfill
\begin{minipage}[t]{.45\textwidth}
{\small 
\begin{verbatim}
    int sum=0;
    for(int i=0;i<n;i++) {
        sum = sum + 1;
    }
    for(int i=0;i<m;i++) {
        sum = sum + 1;
    }
\end{verbatim}
}
\end{minipage}
\begin{sol}
In the code on the left, the inner loop executes $m$ times, so every time
the inner loop executes, $sum$ gets $m$ added to it.  The outer loop executes $n$
times, each time calling the inner loop.  Therefore $m$ is added to $sum$ $n$ times,
so $sum=n\times m$ at the end.

In the code on the right,  The first loop adds $n$ to $sum$, and then the second loop
adds $m$ to $sum$.  Therefore, $sum=n+m$ at the end.
\end{sol}
\end{exa}

The following problem can be solved using the product rule--you just need to figure out how.

\begin{exer}
The number $3$ can be expressed as a sum of one or more positive
integers in four ways, namely, as $3$, $1+2$, $2+1$, and $1+1+1$.
Show that any positive integer $n$ can be so expressed in $2^{n-1}$
ways.

\noindent
\answerlinethree
\begin{answer} 
Write $n = \underbrace{1 + 1 + \cdots + 1}_{n-1 \ +\mathrm{'s}}$.
There are two choices for each plus sign--leave it or perform the addition.
Each of the $2^{n-1}$ ways of making choices leads to a different
expression, and every expression can be constructed this way.
Therefore, there are $2^{n-1}$ such ways of expressing $n$.
\end{answer}
\end{exer}

\begin{exa}
Each day I need to decide between wearing a t-shirt or a polo shirt.
I have 50 t-shirts and 5 polo shirts.  I also have to decide whether to 
wear jeans, shorts, or slacks.  I have 5 pairs of jeans, 15 pairs of shorts,
and 4 pairs of slacks.  How many different choices do I have when I am 
getting dressed?
\begin{sol}
I have $50+5=55$ choices for a shirt and $5+15+4=24$ choices or pants.
So the total number of choices if $55\cdot 24=1320$.
\end{sol}
\end{exa}

\begin{exer}
If license plates are required to have 5 characters, each of which is
either a digits or a letter,
how many license plates are possible?

\noindent
\answerline
\begin{answer}
This combines the product and sum rules.
We now have $10+26=36$ choices for each character, and there are $5$ characters,
so the answer is $36^5$.
\end{answer}
\end{exer}

\begin{exer}
How many bit strings are there of length $n$?

\noindent
\answerline
\begin{answer}
Each bit can be either $0$ or $1$, so there are $2^n$ bit strings of length $n$.
\end{answer}
\end{exer}

\begin{exa}
The integers from $1$ to $1000$ are written in succession. Find the
sum of all the digits.
\begin{sol} When writing the integers
from $000$ to $999$ (with three digits), $3\times 1000 = 3000$
digits are used. Each of the $10$ digits is used an equal number of
times, so each digit is used $300$ times. The the sum of the digits
in the interval $000$ to $999$ is thus
$$(0 + 1 + 2 + 3+ 4 + 5+ 6+ 7 + 8 + 9)\cdot 300 = 13500.
$$Therefore, the sum of the digits when writing the integers from
$1$ to $1000$ is $13500 + 1 = 13501$.
\end{sol}
\end{exa} 

\begin{fib}
In C++, identifiers (e.g. variable and function names) 
can contain only letters (upper and/or lower case), digits, and
the underscore character.  They may not begin with a digit.\footnote{There are 
$84$ reserved keywords that cannot be used, but we will ignore these for this exercise.}
\begin{lenum}
\item There are $26+26+1=53$ possible identifiers that contain a single character.
\item There are $53\cdot(26+26+10+1)=53\cdot 63 = 3339$ 
possible identifiers with two characters.
\item There are \blankline{1.5} possible identifiers that contain three characters.
\item There are \blankline{1.5} possible identifiers that contain four characters.
\item There are \blankline{1.5} possible identifiers that contain $k$ characters.
\end{lenum}
\begin{answer}
$53\cdot 63^2$;
$53\cdot 63^3$;
$53\cdot 63^{k-1}$.
\end{answer}
\end{fib}

\newpage

%%%%%%%%%%%%%%%%%%%%%%%%%%%%%%%%%%%%%%%%%%%%%%%%%%%%%%%%%%%%%%%%%%%%%%
\section{Pigeonhole Principle}\label{sec:pigeon}
The following theorem seems so obvious that it doesn't need to be stated.
However, it often comes in handy in unexpected situations.
\begin{thm}[The Pigeonhole Principle\index{pigeonhole principle}]
If $n$ is a positive integer and $n+1$ or more objects are placed
into $n$ boxes, then one of the boxes contains at least two objects.
\end{thm}

Notice that the pigeonhole principle is saying that this is true
{\em no matter how the objects are places in the boxes}.  In other words,
don't assume that $n-1$ boxes have one object and $1$ box has $2$ objects.
It is possible that all $n+1$ objects are in the same box.  
But no matter how the objects are distributed in the boxes, 
we can be sure that there is some box with at least two objects.

\begin{exa}
In any group of $13$ people, there are always two who have their birthday on the same month.
Similarly, if there are $32$ people, at least two people were born on the same day of the month.
\end{exa}

\begin{exer}
What can you say about the digits in a number that is 11 digits long?

\noindent
\answerline
\begin{answer}
It contains at least one repeated digit.  The wording of your answer is very
important.  Your answer should not be ``it has some digit twice'' since this
is vague--do you mean `exactly twice'?  If so, that is incorrect.  If you mean
`at least twice', then it is better to be explicit and say it that way or just say
`repeated'.  To be clear, we don't know that it contains any digit exactly twice,
and we also don't know how many unique digits the number has--it might be $22222222222$,
but it also might be $98765432101$.
\end{answer}
\end{exer}

The pigeonhole principle can be generalized.
\begin{thm}[The Generalized Pigeonhole Principle]
If $n$ objects are placed into $k$ boxes, then there is at least one box that
contains at least $\lceil n/k\rceil$ objects.
\begin{pf}
Assume not.  Then each of the $k$ boxes contains no more than $\lceil n/k\rceil -1$ objects.
Notice that $\lceil n/k\rceil < n/k+1$ (convince yourself that this is always true).
Thus, the total number of objects in the $k$ boxes is at most
$$k\left(\lceil n/k\rceil -1\right)<k(n/k+1-1)=n,$$
contradicting the fact that there are $n$ objects in the boxes.
Therefore, some box contains at least $\lceil n/k\rceil$ objects.
\end{pf}
\end{thm}

The tricky part about using the pigeonhole principle is identifying 
the {\em boxes} and {\em objects}.  Once that is done, 
applying either form of the pigeonhole principle is straightforward.
Actually, often the trickiest thing is identifying that the pigeonhole
principle even applies to the problem you are trying to solve.

\begin{exa}
A drawer contains an infinite supply of white, black, and blue socks.
What is the smallest number of socks you must take from the drawer in 
order to be guaranteed that you have a matching pair?
\begin{sol}
Clearly I could grab one of each color, so three is not enough.
But according the the Pigeonhole Principle, if I take $4$ socks, then 
I will get at least $\lceil4/3\rceil=2$ of the same color
(the colors correspond to the boxes).  
So $4$ socks will guarantee a matched pair.

Notice that I showed two things in this proof.  I showed that $4$ socks
was enough, but I also showed that $3$ was not enough.  This is important.
For instance, $5$ is enough, but it isn't the smallest number that works.
\end{sol}
\end{exa}

\begin{exer}
An urn contains $28$ blue marbles, $20$ red marbles, $12$ white
marbles, $10$ yellow marbles, and $8$ magenta marbles. How many
marbles must be drawn from the urn in order to assure that there
will be $15$ marbles of the same color? Justify your answer.

\noindent
\answerlinethree
\begin{answer} 
If all
the magenta, all the yellow, all the white, $14$ of the red and $14$
of the blue marbles are drawn, then in among these $8 + 10 + 12 + 14
+ 14 = 58$ there are no $15$ marbles of the same color. Thus we need
$59$ marbles in order to insure that there will be $15$ marbles of
the same color.
\end{answer}
\end{exer}


\begin{exer}
You are in line to get tickets to a concert.
Each person can get at most $4$ tickets.
There are only $100$ tickets available. 
The girl behind you in line says 
``I sure hope there are enough tickets for me.  You're lucky, though.  You will
get as many as you want.''  What does she know, and under what circumstances will
she get any tickets?

\noindent
\answerlinethree
\begin{answer}
She knows that you are the $25$th person in line.  If everyone gets $4$ tickets,
she will get none, but you will get the $4$ you want.  She can get one or more
tickets if one or more people in front of her, including you, get less than $4$.
\end{answer}
\end{exer}

The pigeonhole principle is useful in {\em existence} proofs--that is, 
proofs that show that something exists without actually identifying it concretely.

\begin{exa} 
Show that amongst any seven distinct positive integers
not exceeding $126$, one can find two of them, say $a$ and $b$,
which satisfy $$ b < a \leq 2b.$$
\begin{sol} Split the
numbers $\{ 1, 2, 3, \ldots , 126 \}$ into the six sets $$\{ 1,
2\} , \{ 3, 4, 5, 6\} , \{ 7, 8, \ldots , 13, 14\} , \{ 15, 16,
\ldots , 29, 30\}, $$ $$\{ 31, 32, \ldots , 61, 62\} \ {\rm and} \
\{ 63, 64, \ldots ,126\} .$$By the Pigeonhole Principle, two of
the seven numbers must lie in one of the six sets, and obviously,
any such two will satisfy the stated inequality.
\end{sol}
\end{exa} 

\newpage


\begin{exa}
Given any $9$ integers whose prime factors lie in the set $\{3, 7,
11 \}$ prove that there must be two whose product is a square.
\begin{sol} 
For an integer to be a square, all the exponents of its
prime factorisation must be even. Any integer in the given set has a
prime factorisation of the form $3^a7^b11^c$. Now each triplet $(a,
b, c)$ has one of the following 8 parity patterns: (even, even,
even), (even, even, odd), (even, odd, even), (even, odd, odd), (odd,
even, even), (odd, even, odd), (odd, odd, even), (odd, odd, odd). In
a group of $9$ such integers, there must be two with the same parity
patterns in the exponents. Take these two. Their product is a
square, since the sum of each corresponding exponent will be even.
\end{sol}
\end{exa}

\begin{exer}
The nine entries of a $3\times 3$ grid are filled with $-1$, $0$, or
$1$. Prove that among the eight resulting sums (three columns, three
rows, or two diagonals) there will always be two that add to the
same number.

\noindent
\answerlinethree
\begin{answer}
There are seven possible sums, each one a number in
$\{-3,-2,-1,0,1,2,3\}$. By the Pigeonhole Principle, two of the
eight sums must add up to the same number.
\end{answer}
\end{exer}

\begin{exa}\label{exa:unit_square}
Prove that if five points are taken on or inside a unit square,
there must always be two whose distance is no more than $\dfrac{\sqrt{2}}{2}$.

\noindent
\begin{minipage}[t]{.85\textwidth}
\begin{sol} 
Split the square into four congruent squares as shown to the right. 
At least two of the points must fall into one
of the smaller squares. 
The longest distance between two points in one of the smaller squares is, 
by the
Pythagorean Theorem, $\sqrt{(\frac{1}{2})^2 + (\frac{1}{2})^2} =
\dfrac{\sqrt{2}}{2}$.  Thus, the result holds.
\end{sol}
\end{minipage}
\hfill
\begin{minipage}[t]{.15\textwidth}
$$ \psset{unit=1pc} \psline(-1,-1)(1,-1)(1,1)(-1,1)(-1,-1)
\psline(-1,0)(1,0)
\psline(0,-1)(0,1) $$
\end{minipage}
\end{exa}

\begin{exa} 
Given any set of ten natural numbers between $1$ and $99$
inclusive, prove that there are two distinct nonempty subsets of
the set with equal sums of their elements.
(Hint:  How many possible subsets are there, and what are the possible sums
of the elements within the subsets?)
\begin{sol}
There are $2^{10} - 1 = 1023$ non-empty subsets that one can form
with a given 10-element set. To each of these subsets we associate
the sum of its elements. The minimum value that the sum can be
for any subset is $1+2+\cdots+10=55$, and 
the maximum value is $90 + 91 + \cdots + 99 = 945.$ Since
the number of possible sums is no more than $945-55+1=891<1023$, there
must be at least two different subsets that have the same sum.
\end{sol}
 \end{exa} 
 
%\begin{exa} Prove that if $55$ of the integers from $1$ to $100$ are selected,
%then two of them differ by $10$.
%\begin{sol} First observe that if we choose $n + 1$ integers from
%any set of $2n$ consecutive integers, there will always be some
%two that differ by $n$. This is because we can pair the $2n$
%consecutive integers
%$$ \{ a + 1, a + 2, a + 3, \ldots , a + 2n\}$$ into the $n$ pairs
%$$ \{ a + 1, a + n + 1\}, \{ a + 2, a + n + 2\}, \ldots , \{ a + n, a + 2n\},$$and
%if $n + 1$ integers are chosen from this, there must be two that
%belong to the same group.
%
%So now group the one hundred integers as follows:
%$$ \{ 1, 2, \ldots 20 \} , \{ 21, 22, \ldots , 40\} ,$$  $$ \{ 41, 42, \ldots , 60\}, \
% \{ 61, 62, \ldots , 80\} $$
%and $$ \{ 81, 82, \ldots , 100\}.$$ 
%If we select fifty five integers, then we must have 
%selected at least $\lceil 55/5\rceil = 11$ from one of the groups. 
%From that group, by the above observation (let $n = 10$), there must be
%two that differ by $10$.
%\end{sol}
%\end{exa}

\begin{exer}
An eccentric old man has five cats. These cats have $16$ kittens among
themselves. What is the largest integer $n$ for which one can say
that at least one of the five cats has $n$ kittens?\\
\answerlinetwo
\begin{answer}
We have $ \lceil{\frac{16}{5}}\rceil = 4$, so some cat has at least four kittens.
\end{answer}
\end{exer}

\begin{eval}\label{eval:partyhands}
Prove that at a party with at least two people, 
there are two people who have shaken hands with the same 
number of people.
\begin{spfenum}
\item
There are $n-1$ people 1 person can shake hands with--4 others if there are 5 people at the party.  
At one given time two people cannot shake hands with 0 people and $n-1$ people simultaneously because 
there are 4 slots to fill and 5 people therefore by the pigeonhole principle at least two people 
shake hands with the same number of others.

\noindent
\evallinetwo
\item
Assume that at a gathering of $n\geq 2$ people, there are no two people who have shaken hands with the 
same number of people.  If there are two people at the gathering they must either shake hands with each other 
or shake hands with nobody.  However, this contradicts the assumption that no two people have shaken hands
with the same number of people.  Therefore, by contradiction, at a gathering of $n\geq 2$ people, 
there are at 
least two people who have shaken hands with the same number of people.

\noindent
\evallinetwo
\item
Assume that at a gathering of $n\geq 2$ people, there are no two people who have shaken hands with the 
same number of people.  Person $n$ shakes hands with $n-1$ people because you can't shake your own hand.
Person $n-1$ then shakes hands with $n-2$ people and so on...until you reach the last person.
He shakes hands with no one which fulfills the contradiction.

\noindent
\evallinetwo
\end{spfenum}
\begin{answer}
\begin{pfevalenum}
\item
This proof is incomplete.  It kind of argues it for $5$, not $n$ in general.
Even then, the proof is neither clear not complete.  For instance, what are
the 4 `slots'?
\item
They only prove it for $n=2$.  It needs to be proven for {\em any} $n$.
\item
You can't assume somebody had shaken hands with everyone else without some justification.  
You certainly can't assume it was any particular person (i.e. person $n$).  
Similarly, you can't assume the next person has shaken $n-2$ hands without justifying it. 
The final statement is weird (what does `fulfills the contradiction' mean?)
and needs justification (why is it a problem that the last person shakes no hands?).
\end{pfevalenum}
\end{answer}
\end{eval}

\begin{exer}
Give a correct proof of the problem stated in Evaluate~\ref{eval:partyhands}.
\vspace*{3.25in}
\begin{answer}
Notice that if someone shakes $n-1$ hands, then nobody shakes $0$ hands and vice-verse.
Thus, we have two cases.  
If someone shakes $n-1$ hands, then the $n$ people can shake
hands with between $1$ and $n-1$ other people. 
If nobody shakes hands with $n-1$ people, then the $n$ people can shake
hands with between $0$ and $n-2$ other people. In either case, there are $n-1$
possibilities for the number of hands that the $n$ people can shake.
The pigeonhole principle implies that two people shake hands with the same number of people. 

Note:  You cannot say that the two cases are that someone shakes hands with $n-1$ or
someone shakes hands with $0$. It may be that {\em neither} of these is true.
The two cases are someone shakes hands with $n-1$ others or nobody does.  Alternatively,
you could say someone shakes hands with $0$ others or nobody does.
\end{answer}
\end{exer}

\begin{exer}
There are seventeen friends from high school that all keep in touch by writing letters
to each other.\footnote{You do know what letters are, right?  They are like e-mail, only
they are written on paper, are sent to just one person, and are delivered to your physical mail box.}
%footnote causes underfull hbox, but I don't know how to prevent it.
To be clear, each person writes separate letters to each of the others.
In their letters only three different topics are discussed. 
Each pair only corresponds about one of these topics. 
Prove that there at least three people who all write to each other about the same topic.
\vspace*{3.5in}
\begin{answer} 
Choose a particular person of the group, say Charlie. He
corresponds with sixteen others. By the pigeonhole principle,
Charlie must write to at least six of the people about one topic, say
topic I. If any pair of these six people corresponds about topic I,
then Charlie and this pair do the trick, and we are done.
Otherwise, these six correspond amongst themselves only on topics
II or III. Choose a particular person from this group of six, say
Eric. By the Pigeonhole Principle, there must be three of the five
remaining that correspond with Eric about one of the topics, say
topic II. If amongst these three there is a pair that corresponds
with each other on topic II, then Eric and this pair correspond on
topic II, and we are done. Otherwise, these three people only
correspond with one another on topic III, and we are done again.
\end{answer}
\end{exer}

\newpage

\section{Permutations and Combinations}\label{sec:permComb}

Most of the counting problems we will be dealing with can be classified
into one of four categories. The categories are determined by two factors:
whether or not repetition is allowed and whether or not order matters.
After presenting a brief example of each of these
categories, we will go into more detail about each in the following
four subsections.

\begin{exa}
Consider the set $\{a, b, c, d\}$. Suppose we ``select'' two
letters from these four. Depending on our interpretation, we may
obtain the following answers.
\begin{lenum}
\item {\bf Permutations with repetitions.} The {\em order} of listing the letters is important, and
{\em repetition is} allowed. In this case there are $4\cdot 4 =
16$ possible selections: $$\begin{array}{|l|l|l|l|}\hline aa & ab
& ac & ad
\\ \hline ba & bb & bc & bd \\ \hline ca & cb & cc & cd \\ \hline da & db & dc & dd \\ \hline\end{array}
$$
\item {\bf Permutations without repetitions.} The {\em order} of listing the letters is important, and
{\em repetition is not} allowed. In this case there are $4\cdot 3
= 12$ possible selections: $$\begin{array}{|l|l|l|l|}\hline  & ab
& ac & ad
\\ \hline ba &  & bc & bd \\ \hline ca & cb &  & cd \\ \hline da & db & dc & \\ \hline \end{array}
$$
\item {\bf Combinations with repetitions.}  The {\em order} of listing the letters is {\bf not} important, and
{\em repetition is} allowed. In this case there are
$\dfrac{4\cdot 3}{2}+4 = 10$ possible selections:
$$\begin{array}{|l|l|l|l|} \hline aa & ab & ac & ad
\\ \hline & bb & bc & bd \\ \hline  &  & cc & cd \\ \hline &  &  & dd \\ \hline \end{array}
$$
\item {\bf Combinations without repetitions.}  The {\em order} of listing the letters is {\bf not} important, and
{\em repetition is not} allowed. In this case there are
$\dfrac{4\cdot 3}{2} = 6$ possible selections:
$$\begin{array}{|l|l|l|l|}\hline \hspace{.35cm} & ab & ac & ad
\\ \hline  \hspace*{2ex} &  & bc & bd \\ \hline &  &  & cd \\ \hline  &  &  &  \\ \hline \end{array}
$$
\vspace*{0in}
\end{lenum}
\end{exa}

Although most of the simple types of counting problems we want to solve can be reduced to 
one of these four, care must be taken.  The previous example assumed that we had a set of
{\em distinguishable} objects.  When objects are not distinguishable, the situation is more complicated.  

\subsection{Permutations without Repetitions}

\begin{df}\index{permutation}
Let $x_1, x_2, \ldots , x_n$ be $n$ distinct objects. A {\bf permutation} 
of these objects is simply a rearrangement of them.
\end{df}
\begin{exa}
There are $24$ permutations of the letters in $MATH$, namely
$$\begin{array}{llllll} MATH & MAHT & MTAH & MTHA & MHTA &  MHAT\\
 AMTH & AMHT & ATMH & ATHM & AHTM &  AHMT\\
  TAMH & TAHM & TMAH & TMHA & THMA &  THAM\\
   HATM & HAMT & HTAM & HTMA & HMTA &  HMAT\\       
\end{array}   
$$
\end{exa}

\begin{exer}
List all of the permutations of $EAT$

\noindent
\answerlinetwo
\begin{answer}
$EAT$, $ETA$, $ATE$, $AET$, $TAE$, and $TEA$.
\end{answer}
\end{exer}


\begin{thm}\label{thm:counting_permutations}
Let $x_1, x_2, \ldots , x_n$ be $n$ distinct objects. Then there are
$n!$ permutations of them.
\begin{pf}
The first position can be chosen in $n$ ways, the second object in
$n - 1$ ways, the third in $n - 2$, etc. This gives $$n(n-1)(n-2)
\cdots 2\cdot 1 = n!.
$$
\end{pf}
\end{thm}

\begin{exa}
Previously we saw that there are $24=4!$ permutations of the letters in $MATH$ and
$6=3!$ permutations of the letters in $EAT$.
\end{exa}

\begin{exer}
How many permutations are there of the letters in {\sc uncopyrightable}?

\noindent
\answerline
\begin{answer}
Since there are $15$ letters and none of them repeat, there are $15!$ permutations
of the letters in the word {\sc uncopyrightable}.
\end{answer}
\end{exer}

% Consider adding this back as some point, but perhaps a different algorithm, and
%certainly with more explanation.
% CAC, 6/17/14
% Not correct as written--the algorithm only gives the next permutation.
% Also, would need more explanation to be useful.
%CAC, 8/7/13
%\begin{exa}
%Write a code fragment that prints all $n!$ of the set $\{1,2,\ldots
%, n\}$.
%\end{exa}\begin{sol} The following algorithm prints them in
%lexicographical order. We use the algorithms from
%Examples \ref{exa:swapping1} and \ref{exa:ReverseArray}.
%
% \bcp[Ovalbox]{Permutations}{n}
%k \GETS n -1 \\
%\WHILE X[k]>X[k-1] \DO \BEGIN k \GETS k-1 \END \\
%t \GETS k+1 \\
%\WHILE ((t<n) \AND (X[t+1]>X[k])) \DO \BEGIN t\GETS t+1 \END \\
%\COMMENT \rm{now\ }X[k+1]>\ldots > X[t]>X[k]>X[t+1]>\ldots > X[n]\\
%\sc{Swap}(X[k], X[t]) \\
%\COMMENT \rm{now\ }X[k+1]>\ldots > X[n]\\
%\sc{ReverseArray}(X[k+1], \ldots , X[n])
% \ecp
%\end{sol}

Let's see some slightly more complicated examples.

\begin{exa}
A bookshelf contains $5$ German books, $7$ Spanish books and $8$
French books. Each book is different from one another.
How many different arrangements can be done of these books if
\begin{lenum}
\item  we put no restrictions on how they can be arranged?
\item  books of each language must be next to each other? 
\item all the French books must be next to each other? 
\item no two French books must be next to each other?
\end{lenum}
\begin{sol}
\begin{lenum} \item We are permuting $5 + 7 + 8 = 20$ objects. Thus the number
of arrangements sought is $20! = 2432902008176640000$. \item
``Glue'' the books by language, this will assure that books of the
same language are together. We permute the $3$ languages in $3!$
ways. We permute the German books in $5!$ ways, the Spanish books in
$7!$ ways and the French books in $8!$ ways. Hence the total number
of ways is $3!\cdot 5!\cdot 7!\cdot 8! = 146313216000$. 
\item 
Align the German books
and the Spanish books first. Putting these $5 + 7 = 12$ books
creates $12 + 1 = 13$ spaces (we count the space before the first
book, the spaces between books and the space after the last book).
To assure that all the French books are next each other, we ``glue''
them together and put them in one of these spaces.  Now, the French
books can be permuted in $8!$ ways and the non-French books can be
permuted in $12!$ ways. Thus the total number of permutations is
$$ 13\cdot 8!\cdot 12! =  251073478656000.$$
\item 
As with (c), we align the $12$ German and Spanish books first,
creating $13$ spaces. To assure that no two French books are next to
each other, we put them into these spaces. The first French book can
be put into any of $13$ spaces, the second into any of $12$ remaining spaces, etc.,
and the eighth French book can be put into any $6$ remaining spaces. Now, the
non-French books can be permuted in  $12!$ ways. Thus the total
number of permutations is
$$ 13\cdot 12\cdot 11\cdot 10\cdot 9\cdot 8\cdot 7\cdot 6\cdot 12!=24856274386944000.$$
\end{lenum}
\end{sol}
\end{exa}



\begin{exer}
Telephone numbers in {\em Land of the Flying Camels} have $7$
digits, and the only digits available are $\{0,1, 2,3,4,5, 7, 8\}$.
No telephone number may begin in $0$, $1$ or $5$. Find the number of
telephone numbers possible that meet the following criteria:
\begin{lenum}
 \item 
 You may not repeat any of the digits.
 
\answerlinetwo
\item 
You may not repeat the digits and the phone numbers must be odd.

\answerlinetwo
\end{lenum}
\begin{answer}
(a) $5\cdot 7\cdot 6 \cdot 5 \cdot 4 \cdot 3 \cdot 2 = 25,200$. 
(b) We condition on the last digit. 
If the last digit were $1$ or $5$ then we would have $5$
choices for the first digit and $2$ for the last digit.
Then there are $6$ left to choose from for the second, $5$ for the third, etc.
So this leads to  
$$5\cdot 6 \cdot 5\cdot 4 \cdot 3\cdot 2 \cdot 2 = 7,200$$ 
possible phone numbers. 
If the last digit were either $3$ or $7$, then 
we would have $4$ choices for the first digit and $2$ for the last.
The rest of the digits have the same number of possibilities as above, so we would
have
$$4\cdot 6\cdot 5 \cdot 4 \cdot 3\cdot 2 \cdot 2  = 5,760   $$
possible phone numbers. Thus the total number of phone numbers is 
$$7200 + 5760 = 12,960.$$
\end{answer}
\end{exer}

The previous example and exercise should demonstrate that counting often requires thinking about
things in different ways depending on the exact situation.  This can be tricky, and it
is very easy to make mistakes that lead to under or over counting possibilities.
As you are solving problems, think very carefully about what you are counting so 
you don't fall into this trap.

\subsection{Permutations with Repetitions}
We now consider permutations with repeated objects.

\begin{exa}\label{exa:permu_repetitions}
In how many ways may the letters of the word $MASSACHUSETTS$ be
permuted to form different strings?
\begin{sol} We put subscripts on the repeats forming
$$MA_1S_1S_2A_2CHUS_3ET_1T_2S_4.$$There are now $13$
distinguishable objects, which can be permuted in $13!$ different
ways by Theorem \ref{thm:counting_permutations}. 
But this counts some arrangements multiple times since in reality
the duplicated letters are not distinguishable.
Consider a single permutation of all of the distinguishable letters.
If I permute the letters $A_1A_2$, I get the same permutation when ignoring the subscripts.
The same thing is true of $T_1T_2$. Similarly, there are $4!$ permutations of $S_1S_2S_3S_4$,
so there are $4!$ permutations that look the same (without the subscripts).
Since I can do all of these independently, there are $2!2!4!$ permutations that look identical
when the subscripts are removed.  This is true of every permutation.  
Therefore, the actual number of permutations is 
$\displaystyle\frac{13!}{2!\cdot 4!\cdot 2!} = 64864800.$
\end{sol}
\end{exa}

The following exercises should help the technique used in the previous example to sink in.

\begin{exer}\label{exer:tallperm1}
Use an argument similar to that in Example~\ref{exa:permu_repetitions} to 
determine the number of permutations in the letters in $TALL$.

\noindent
\answerlinetwo
\begin{answer}
Label the letters $T_1$, $A_1$, $L_1$, and $L_2$.  There are $4!$ permutations of these letters.
However, every permutation that has $L_1$ before $L_2$ is actually identical to one having
$L_1$ before $L_2$, so we have double-counted.  Therefore, there are $4!/2=12$ permutations
of the letters in $TALL$.
\end{answer}
\end{exer}

\begin{exer}
List all of the permutations of the letters $TALL$.

\noindent
\answerlinetwo
\begin{answer}
$TALL$, $TLAL$, $TLLA$, $ATLL$, $ALTL$, $ALLT$, $LLAT$, $LALT$, $LATL$, $LLTA$, $LTLA$, and $LTAL$.
That makes $12$ permutations, which is exactly what we said it should be in Exercise~\ref{exer:tallperm1}.
\end{answer}
\end{exer}

\begin{exer}
How many permutations are there in the letters of {\em AEEEI}?

\noindent
\answerlinetwo
\begin{answer}
Following similar logic to the previous few examples, since we have one letter that is repeated
three times, and a total of $5$ letters, the answer is $5!/3!=20$.
\end{answer}
\end{exer}

\begin{exer}
List all of the permutations of the letters $AEEEI$.

\noindent
\answerlinethree
\begin{answer}
Ten of them are
$AIEEE$, $AEIEE$, $AEEIE$, $AEEEI$, $EAIEE$, $EAEIE$, $EAEEI$, $EEAIE$, $EEAEI$, $EEEAI$.
The other ten are identical to these, but with the $A$ and $I$ swapped.
\end{answer}
\end{exer}

The arguments from the previous examples and exercises can be generalized to prove the following.

\begin{thm}\label{thm:genPerms}
Let there be $k$ types of objects: $n_1$ of type $1$;  $n_2$ of type
$2$; etc. Then the number of ways in which these $n_1 + n_2 + \cdots
+n_k$ objects can be rearranged is $$ \frac{(n_1 + n_2 + \cdots
+n_k)!}{n_1!\cdot n_2! \cdots n_k!}. $$
\end{thm}


\begin{exa} How many permutations of the letters from
$MASSACHUSETTS$ contain $MASS$?
\begin{sol} 
We can consider $MASS$ as one block along with the remaining
$9$ letters $A,$ $C,$ $H,$ $U,$ $S,$ $E,$ $T,$ $T,$ $S$. Thus, we are
permuting $10$ `letters'.
There are two $S$'s\footnote{Remember, the other two $S$'s are part of $MASS$, which we are
now treating as a single object.} 
and two $T$'s and so the total number of permutations sought is
$$\frac{10!}{2!\cdot 2!} = 907200.
$$
\end{sol}
\end{exa}

\begin{exer} 
How many permutations of the letters from the word 
$ALGORITHMS$ contain $SMITH$?

\noindent
\answerlinethree
\begin{answer} 
We can consider $SMITH$ as one block along with the remaining
$5$ letters $A$, $L$, $G$, $O$, and $R$. Thus, we are
permuting $6$ `letters', all of which are unique.  So there are
$6!=720$ possible permutations.
\end{answer}
\end{exer}


\begin{exa}\label{exa:summingto9}
In how many ways may we write the number $9$ as the sum of three
positive integer summands? Here order counts, so, for example, $1 +
7 + 1$ is to be regarded different from $7 + 1 + 1$.
\begin{sol} 
We need to find the values of $a$, $b$, and $c$ such that
$a + b + c = 9,$
where $a,b,c\in\BBZ^+$.
We will consider triples $(a,b,c)$ listed smallest
to largest and determine how many ways each triple can be reordered.
The possibilities are:
$$
\begin{array}{|l|l|}\hline 
(a,b,c) & \mathrm{Number\ of\ permutations} \\
\hline (1,1,7) & 3!/2! = 3\\
\hline (1,2,6) & 3! = 6\\
\hline (1,3,5) & 3! = 6\\
\hline (1,4,4) & 3!/2!  = 3\\
\hline (2,2,5) & 3!/2!  = 3\\
\hline (2,3,4) & 3! = 6\\
\hline (3,3,3) & 3!/3!  = 1\\
\hline
\end{array}$$
Thus the number desired is $3 + 6 + 6 +3 + 3 + 6 + 1 = 28.$
\end{sol}
\end{exa}

\begin{exa}
In how many ways can the letters of the word {\bf MURMUR} be
arranged without allowing two of the same letters next to each other?
\begin{sol} If we started with, say , {\bf MU} then the {\bf R} could
be arranged in one of the following three ways:
$$\begin{array}{|c|c|c|c|c|c|}\hline 
{\bf M} & {\bf U} & {\bf R} & \hspace{2ex}  & {\bf R} &  \hspace{2ex}\\ 
\hline
\multicolumn{6}{c}{}\\
\hline 
{\bf M} & {\bf U} & {\bf R} &  \hspace{2ex}   & \hspace{2ex} & {\bf R}\\ 
\hline
\multicolumn{6}{c}{}\\
\hline 
{\bf M} & {\bf U} &  \hspace{2ex}  & {\bf R}  & \hspace{2ex} & {\bf R}\\ 
\hline 
\end{array}
$$
In the first case there are $2! = 2$ ways of putting the remaining {\bf
M} and {\bf U}, in the second there are $2!=2$ ways and in the third
there is only $1!$ way. Thus starting the word with {\bf MU} gives $2 +
2 + 1 = 5$ possible arrangements. In the general case, we can choose
the first letter of the word in $3$ ways, and the second in $2$
ways. Thus the number of ways sought is $3\cdot 2 \cdot 5 = 30$.\footnote{It should
be noted that this analysis worked because the three letters each occurred twice.
If this was not the case we would have had to work harder to solve the problem.}
\end{sol}
\end{exa}

\newpage

\begin{exer}
Telephone numbers in {\em Land of the Flying Camels} have $7$
digits, and the only digits available are $\{0,1, 2,3,4,5, 7, 8\}$.
No telephone number may begin with $0$, $1$ or $5$. Find the number of
telephone numbers possible that meet the following criteria:
\begin{lenum}
\item You may repeat all digits.

\answerlinetwo
 \item You may repeat digits, but the last digit must be even.
 
\answerlinetwo
\item You may repeat digits, but the last digit must be odd.

\answerlinetwo
\end{lenum}
\begin{answer}
(a) $5\cdot 8^6 =1310720$.
(b) $5\cdot 8^5 \cdot 4 = 655360 $.
(c) $5\cdot 8^5 \cdot 4 = 655360 $.
\end{answer}
\end{exer}


\begin{exa}
In how many ways can the letters of the word {\bf AFFECTION} be
arranged, keeping the vowels in their natural order and not letting
the two {\bf F}'s come together?
\begin{sol} There are $\dfrac{9!}{2!}$ ways of permuting the letters
of {\bf AFFECTION}. The $4$ vowels can be permuted in  $4!$ ways,
and in only one of these will they be in their natural order. Thus
there are $\dfrac{9!}{2!\cdot 4!}$ ways of permuting the letters of {\bf
AFFECTION} in which their vowels keep their natural order.
If we treat $FF$ as a single letter, there are $8!$ ways of permuting
the letters so that the $F$'s stay together.  Hence there
are $\dfrac{8!}{4!}$ permutations of {\bf AFFECTION} where
the vowels occur in their natural order and the $FF$'s are together.
In conclusion, the number of permutations sought is
$$ \dfrac{9!}{2!\cdot 4!} - \dfrac{8!}{4!} = \dfrac{8!}{4!}\left(\dfrac{9}{2} -1\right) 
= 8\cdot 7\cdot 6 \cdot 5\cdot\dfrac{7}{2}  = 5880.$$
\end{sol}
\end{exa}

\newpage

\subsection{Combinations without Repetitions}
Let's begin with some important notation.

\begin{df}\index{binomial coefficient}\index{choose}\index{a choose@$\binom{n}{k}$ (binomial coefficient)}
Let $n, k$ be  non-negative integers with $0\leq k \leq n$. The
{\bf binomial coefficient} $\dis{\binom{n}{k}}$ (read ``$n$ {\em choose} $k$'') is
defined by
$$\binom{n}{k} = \frac{n!}{k!(n - k)!} = \frac{n\cdot (n - 1) \cdot (n - 2) \cdots (n - k + 1)}{1\cdot 2 \cdot 3 \cdots k}.$$
An alternative notation is $C(n,k)$.  This notation is particularly useful when you want to express
a binomial coefficient in the middle of text since it doesn't take up two lines.
\end{df}

\begin{rem}
Observe that in the last fraction, there are $k$ factors in both the
numerator and denominator. Also, observe the boundary conditions
$$\binom{n}{0} = \binom{n}{n} = 1,\ \ \ \ \binom{n}{1} =
\binom{n}{n-1} = n.
$$
\end{rem}

\begin{exa}We have
\begin{eqnarray*}
\binom{6}{3} & =  &\frac{6\cdot 5 \cdot 4}{1\cdot 2 \cdot 3} = 20,  \\
 \binom{11}{2}  & =  & \frac{11\cdot 10}{1\cdot 2 } = 55,  \\ \binom{12}{7} & =  &
 \frac{12\cdot 11 \cdot 10\cdot 9 \cdot 8 \cdot 7
\cdot 6}{1\cdot 2 \cdot 3\cdot 4\cdot 5 \cdot 6 \cdot 7} = 792,  \\
\binom{110}{0} &   =  & 1.
\end{eqnarray*}
\end{exa}

\begin{exer}
Compute each of the following
\begin{lenum}
\item
$\displaystyle\binom{7}{5} =\blankline{2}$
\item
 $\displaystyle\binom{12}{2}  = \blankline{2}$
\item 
 $\displaystyle\binom{10}{5} =\blankline{2}$
\item
$\displaystyle\binom{200}{4} = \blankline{2}$
\item
$\displaystyle\binom{67}{0} = \blankline{2}$
\end{lenum}
\begin{answer}
(a) $\dfrac{7\cdot 6\cdot 5 \cdot 4\cdot 3}{1\cdot 2\cdot 3\cdot 4 \cdot 5} = 21$.
(b) $\dfrac{12\cdot 11}{1\cdot 2} = 66$.
(c) $\dfrac{10\cdot 9\cdot 8\cdot 7\cdot 6}{1\cdot 2\cdot 3\cdot 4 \cdot 5} = 252$.

\noindent
(d) $\dfrac{200\cdot 199\cdot 198\cdot 197}{1\cdot 2\cdot 3\cdot 4} =64,684,950$.
(e) $1$.
\end{answer}
\end{exer}

If there are $n$ kittens and you decide to take $k$ of them home, you also
decided {\em not} to take $n-k$ of them home.  This idea leads to the following 
important theorem.
\begin{thm}\label{thm:binomminus}
If $n, k\in \BBZ$, with $0 \leq k\leq n$, then
$$\dis{\binom{n}{k} = \frac{n!}{k!(n - k)!} = \frac{n!}{(n-k)!(n -
(n-k))!} = \binom{n}{n - k}}$$
\begin{pf}
Since $k=n-(n-k)$, the result is obvious.
\end{pf}
\end{thm}

\begin{exa}
$$\binom{11}{9} = \binom{11}{2} = 55.$$
$$\binom{12}{5} = \binom{12}{7} = 792.$$
$$\binom{110}{109} = \binom{110}{1}  =  110$$ 
\end{exa}


\begin{exer}
Compute each of the following
\begin{lenum}
\item
$\displaystyle\binom{17}{15} =\blankline{2}$
\item
 $\displaystyle\binom{12}{10}  = \blankline{2}$
%\item 
% $\displaystyle\binom{10}{6} =\blankline{2}$
\item
$\displaystyle\binom{200}{196} = \blankline{2}$
\item
$\displaystyle\binom{67}{66} = \blankline{2}$
\end{lenum}
\begin{answer}
(a) $\displaystyle\binom{17}{15} = \displaystyle\binom{17}{2}=\dfrac{17\cdot 16}{1\cdot 2} = 136$.
(b) $\displaystyle\binom{12}{10} = \displaystyle\binom{12}{2}=\dfrac{12\cdot 11}{1\cdot 2} = 66$.

\noindent
%(c) $\displaystyle\binom{10}{6}=\displaystyle\binom{10}{4}=\dfrac{10\cdot 9\cdot 8\cdot 7}{1\cdot 2\cdot 3\cdot 4} = 210$.
(c) $\displaystyle\binom{200}{196}=\displaystyle\binom{200}{4}=\dfrac{200\cdot 199\cdot 198\cdot 197}{1\cdot 2\cdot 3\cdot 4} =64,684,950$.
%\noindent
(d) $\displaystyle\binom{67}{66}=\displaystyle\binom{67}{1}=67/1=67$.
\end{answer}
\end{exer}

\newpage

\begin{df}\index{combination}
Let there be $n$ distinguishable objects. A $k$-{\bf combination} is
a selection of $k$, ($0 \leq k \leq n$) objects from the $n$ made
without regards to order.
\end{df}
\begin{exa} The
2-combinations from the list $\{X, Y, Z, W\}$ are
$$XY, XZ, XW, YZ, YW, WZ.$$

\noindent
Notice that $YX$ (for instance) is not on the list because $XY$ is already
on the list and order does not matter.
\end{exa}

\begin{exa} The
3-combinations from the list $\{X, Y, Z,W\}$ are
$$XYZ, XYW, XZW,  YWZ.$$
\end{exa}


\begin{exer} List the 2-combinations from the list $\{1, 2, 3, 4, 5\}$

\noindent
\answerlinetwo
\begin{answer}
$12, 13, 14, 15, 23, 24, 25, 34, 35, 45$.
\end{answer}
\end{exer}

\begin{thm}
Let there be $n$ distinguishable objects, and let $k$, $0 \leq k
\leq n$. Then the numbers of $k$-combinations of these $n$ objects
is $\dis{\binom{n}{k}}$.
\begin{pf}
The number of ways of picking $k$ objects if the order matters is 
$n(n - 1)(n - 2) \cdots (n - k + 1)$ since there are $n$ ways of choosing the 
first object, $n - 1$ ways of choosing the second object, etc..
Since each $k$-combination can be ordered in $k!$ ways, the number
of ordered lists of size $k$ is $k!$ times the number of $k$-combinations.
Put another way, the number of $k$-combinations is the number above divided by $k!$.  
That is, the total number of $k$-combinations is 
$$ \frac{n(n - 1)(n - 2) \cdots (n - k + 1)}{k!}=  \binom{n}{k}.$$
\end{pf}
\end{thm}

\begin{exa}
From a group of $10$ people, we may choose a committee of $4$ in
$\dis{\binom{10}{4} = 210} $ ways.
\end{exa}

\newpage
\begin{eval}
A family has seven women and nine men.  They need five of them to get together to plan a party.
If at least one of the five must be a woman, how many ways are there to select the five?
\begin{ssolenum}
\item
Since one has to be a woman, this is equivalent to selecting four people from a pool of 15,
so the answer is $\binom{15}{4}$.

\noindent
\evallinetwo
\item
There are $7$ women to choose from to ensure there is one woman, and then $4$ more need
to be selected from the remaining $15$.  There are $\binom{15}{4}$ ways of doing that.
Therefore the total number of ways is $\binom{15}{4}+7$.

\noindent
\evallinetwo
\item
There are $\binom{16}{5}$ possible committees, $\binom{9}{5}$ of which contain only men.
Thus, there are $\binom{16}{5}-\binom{9}{5}$ committees that contain at least one woman.

\noindent
\evallinetwo
\end{ssolenum}
\begin{answer}
\begin{solevalenum}
\item
This solution does not take into account which woman was selected 
and which 15 of the original 16 are left, so this is not correct.
\item
This solution has two problems.  First, it counts things multiple times.  For instance,
any selection that contains both Sally and Kim will be counted twice--once
when Sally is the first woman selected and again when Kim is selected first.
Second, the product rule should have been used instead of the sum rule.  Of course, that
hardly matters since it would have been wrong anyway.
\item
This solution is correct.
\end{solevalenum}
\end{answer}
\end{eval}


\begin{exa}\label{exa:routes}
Consider the following grid:

\begin{pspicture}(-2, -0.25)(6,3)
\rput(-2,0){\psset{unit=2pc}\psset{gridwidth=1pt,gridlabels=0, subgriddiv=0}\psgrid(0,0)(0,0)(6,3)\psline[linewidth=2pt, linecolor=red](0,0)(4,0)(4,2)(6,2)(6,3)
\uput[l](0,0){$A$} \uput[r](6,3){$B$} }
\end{pspicture}

To count the  number of shortest routes from $A$ to $B$ (one of which is given), 
observe that any shortest path must consist of $6$
horizontal moves and $3$ vertical ones for a total of $6 + 3 = 9$
moves. 
Once we choose which $6$ of these $9$ moves are horizontal the 
$3$ vertical ones are determined.
For instance, if I choose to go horizontal on moves $1$, $2$, $4$, $6$, $7$, and $8$, 
then moves $3$, $5$ and $9$ must be vertical.
Since there are $9$ moves, I just need to choose which $6$ of these are the horizontal
moves.  Thus there are $\binom{9}{6} = 84$ paths.

Another way to think about it is that we need to compute the number of permutations
of $EEEEEENNN$, where $E$ means move east, and $N$ means move north.  
The number of permutations is $9!/(6!\cdot 3!)=\binom{9}{6}$.
\end{exa}

\newpage
\begin{exer}\label{exa:routes2}
Count the number of shortest routes from $A$ to $B$ that pass through point $O$
in the following grid.  (Hint:  Break it into two subproblems and combine the solutions.)

\begin{pspicture}(-2, -0.25)(6,3)
\rput(-2,0){\psset{unit=2pc}\psset{gridwidth=1pt,gridlabels=0, subgriddiv=0}\psgrid(0,0)(0,0)(6,3)
\psdots[dotstyle=*, dotscale=1](3,2)
\psline[linewidth=2pt, linecolor=red](0,0)(3,0)(3,2)(6,2)(6,3)
\uput[l](0,0){$A$}\uput[ur](3,2){$O$} \uput[r](6,3){$B$}}
\end{pspicture}
\vspace*{1.125in}

\begin{answer}
To count the  number of shortest routes from $A$ to $B$ 
that pass through point $O$, we count the number of
paths from $A$ to $O$ (of which there are $\binom{5}{3} = 10$) and
the number of paths from $O$ to $B$ (of which there are
$\binom{4}{3} = 4$). Using the product rule, 
the desired number of paths is
$\displaystyle\binom{5}{3}\binom{4}{3} = 10\cdot 4 = 40$.
\end{answer}
\end{exer}

\begin{eval}
A family has seven women and nine men.
How many ways are there to select five of them to plan a party 
if at least one man and one woman must be selected?
\begin{ssolenum}
\item
There are $7$ choices for the first woman, $9$ choices for the first man, and
$\binom{14}{3}$ choices for the rest of the committee.  Thus, there are
$\binom{14}{3}\cdot 7 \cdot 9$ possible committees.

\noindent
\evallinetwo
\item
Since one has to be a woman and one has to be a man, then they really just need to select
3 more member from the remaining 14 people, so the answer is $\binom{14}{3}$.

\noindent
\evallinetwo
\end{ssolenum}
\begin{answer}
\begin{solevalenum}
\item
This answer is incorrect since it will count some of the committees multiple times.  
If you did not come up with an example of something that gets counted multiple times,
you should do so to convince yourself that this answer is incorrect.
\item
This solution is incorrect since it does not take into account which man and woman 
were selected  and which 14 of the original 16 are left.
\end{solevalenum}
\end{answer}
\end{eval}

Now it's your turn to give a correct solution to the previous problem.
\begin{exer}
A family has seven women and nine men.
How many ways are there to select five of them to plan a party 
if at least one man and one woman must be selected?

\vspace*{1.1in}

\begin{answer}
There are $\binom{16}{5}$ possible committees.  Of these, $\binom{9}{5}$  contain only men and 
$\binom{7}{5}$ contain only women.  
Clearly these two sets of committees do not overlap.
Therefore, the number of committees that contain at least one man and at least one woman is
$\binom{16}{5}-\binom{9}{5} -\binom{7}{5}$.
\end{answer}
\end{exer}

\begin{ques}
In the answer to the previous problem, we pointed out that two sets of committees did not
overlap.  Why was that important?

\noindent
\answerlinetwo
\begin{answer}
Because we subtracted the size of both of these from the total number of possible
committees.  If the sets intersected, we would have subtracted some possibilities twice 
and the answer would have been incorrect.
\end{answer}
\end{ques}

\begin{exa}
Three different integers are drawn from the set $\{1,2,\ldots ,
20\}$. In how many ways may they be drawn so that their sum is
divisible by $3$?
\begin{sol} In $\{1,2,\ldots , 20\}$ there are
\begin{center}
\begin{tabular}{ll}
$6$ & numbers leaving remainder $0$ \\
$7$ & numbers leaving remainder $1$ \\
$7$ & numbers leaving remainder $2$ \\
\end{tabular}
\end{center}
 The sum
of three numbers will be divisible by $3$ when (a) the three numbers
are divisible by $3$; (b) one of the numbers is divisible by $3$,
one leaves remainder $1$ and the third leaves remainder $2$ upon
division by $3$; (c) all three leave remainder $1$ upon division by
$3$; (d) all three leave remainder $2$ upon division by $3$. Hence
the number of ways is
$$\binom{6}{3} + \binom{6}{1}\binom{7}{1}\binom{7}{1} + \binom{7}{3}+\binom{7}{3} = 384.$$
\end{sol}
\end{exa}



\begin{eval}\label{eval:courses}
The 300-level courses in the CS department are split into three groups: Foundations (361, 385), 
Applications (321, 342, 392), and Systems (335, 354, 376).  
In order to get a BS in computer science at Hope you need to take at least one course from
each group.
If you take four 300-level courses, 
how many different possibilities do you have that satisfy the requirements?
\begin{ssolenum}
\item
You have to take one from each group and then you can take any of the remaining 5 courses.
So the total is $2*3*3*5=90$.

\noindent
\evallinetwo
\item
$\binom{8}{4}=70$

\noindent
\evallinetwo
\end{ssolenum}
\begin{answer}
\begin{solevalenum}
\item
This solution is incorrect since it double counts some of the possibilities.
\item
This solution is incorrect because it does not take into account the requirement
that one course from each group must be taken.
\end{solevalenum}
\end{answer}
\end{eval}


\begin{eval}
Using the same requirements from Evaluate~\ref{eval:courses}, 
how many total ways are there to take 300-level courses that satisfy the requirements?
\begin{ssolenum}
\item
Take one from each group and then choose between 0 and 5 of the remaining 5.  The total is
therefore $2*3*3*\dsum_{k=0}^5 \binom{5}{k}$.

\noindent
\evallinetwo
\item
Since you can take anywhere between $3$ and $8$ courses, the number of possibilities is
$\binom{8}{3}+\binom{8}{4}+\binom{8}{5}+\binom{8}{6}+\binom{8}{7}+\binom{8}{8}$.

\noindent
\evallinetwo
\end{ssolenum}
\begin{answer}
\begin{solevalenum}
\item
This solution is incorrect since it counts some of the possibilities multiple times.
\item
This solution is incorrect because it does not take into account the requirement
that one course from each group must be taken.
\end{solevalenum}
\end{answer}
\end{eval}

In Problem~\ref{prob:csCourses} you will have a chance to properly 
solve the problems from the previous two Evaluate exercises.

\subsection{Combinations with Repetitions}

\begin{exa}
How many ways are there to put $10$ ping pong balls into $4$ buckets?
\begin{sol}
We will solve this using a technique sometimes called {\em bars and stars}.
We will represent the drawers with bars and the balls with stars.  We will use
$10$ stars and $3$ bars.  To see why this is $3$ and not $4$, let's see how we 
represent the situation of having $3$ balls in the first bucket, $5$ in the second,
none in the third, and $2$ in the fourth:
\begin{verbatim}
         ***|*****||**
\end{verbatim}
Do you see it?  The bars act as separators between the buckets, which is why we
have one less bar than the number of buckets.

Given this formulation, aren't we just trying to find all possible orderings of
bars and stars?  Indeed.  To do so, all we need to do is determine where to put the
stars, and the bars `fall into place'.  Alternatively, we can determine where to
put the bars and let the stars fall into place.  
There are $13$ spots and we need to choose $10$ spots for the balls (the `stars')
or $3$ spots for the bucket separators (the `bars').
So the solution is 
$$\binom{13}{10}=\binom{13}{3}=286.$$
Notice that Theorem~\ref{thm:binomminus} implies that these two methods
of solving the problem will always be the same, which is a really good thing.
\end{sol}
\end{exa}

\begin{exa}
How many ways are there to choose $10$ pieces of fruit if you can take any number
of bananas, oranges, apples, or pears and the order I select them does not matter?
\begin{sol}
Again we can use stars and bars so solve this problem.  We need $10$ stars to represent
the chosen fruits and $3$ bars to divide the four fruits we can choose from.
The stars before the first bar represent bananas, those between the first and second
bar are oranges, between the second and third are apples, and after the third are pears.
Thus, we need to count the number of ways we can arrange $10$ stars and $3$ bars.
Notice that this is exactly the same thing we did in the previous example, so the
answer is 
$$\binom{13}{10}=\binom{13}{3}=286.$$
\end{sol}
\end{exa}

\begin{exer}
I want to make a sandwich that has $3$ slices of meat.
My refrigerator is well stocked because I have $11$ different meats to choose from.
How many choices do I have for my sandwich if I allow myself to have
multiple slices of the same meat and the order the slices appear on the sandwich
does not matter?

\noindent
\answerlinethree
\begin{answer}
Using $10$ bars to separate the meat and $3$ stars to represent the slices,
we can see that this is exactly the same as the previous two examples.  Thus, 
the solution is $\binom{13}{10}=\binom{13}{3}=286.$
\end{answer}
\end{exer}

We can generalize the previous examples as follows.
\begin{thm}\label{thm:ballsbins}
There are $\displaystyle\binom{n+k-1}{k}=\binom{n+k-1}{n-1}$ ways of placing
$k$ indistinguishable objects into $n$ distinguishable bins.

This is also the number of ways of selecting $k$ objects from a collection of
$n$ objects if repetition is allowed.
\end{thm}

The previous theorem can be applied to various situations.  As with the pigeonhole
principle, the trickiest part is recognizing when and how to apply it.


%Let's see another sort of counting problem.
%
%\begin{thm}[De Moivre]
%Let $n$ be a positive integer. The number of positive integer
%solutions to
%$$x_1 + x_2 + \cdots + x_r = n$$is
%$$\binom{n - 1}{r - 1}.$$
%\label{thm:distributions}
%\begin{pf}
%Write $n$ as
%$$n = 1 + 1 + \cdots + 1 + 1,$$
%where there are $n$ 1s and $n - 1$
%$+$s. To decompose $n$ in $r$ summands we only need to choose $r -1$ pluses from the $n - 1$.
%For instance, writing $n=7$ as $7=2+3+2$ is equivalent to $7=(1+1)+(1+1+1)+(1+1)$, where the $+$'s
%outside of the parentheses are the ones we chose (we apply the ones inside the parentheses to
%obtain the $x_i$s).
%This proves the theorem.
%\end{pf}
%\end{thm}
%
%\begin{exa}\label{exa:blah}
%In how many ways may we write the number $9$ as the sum of three
%positive integer summands? Here order counts, so, for example, $1 +
%7 + 1$ is to be regarded different from $7 + 1 + 1$.
%\begin{sol} Notice that this is the same problem as Example \ref{exa:summingto9}. We are
%seeking integral solutions to $$ a + b + c = 9,  \ \ \ a>0, b>0, c>0.$$ 
%By Theorem \ref{thm:distributions} this is $$\binom{9-1}{3-1} = \binom{8}{2} =  28.$$
%\end{sol}
%\end{exa}
%
%\begin{rem}
%The solution in Example \ref{exa:blah} was much easier than the solution in Example \ref{exa:summingto9},
%demonstrating the fact that choosing the right tool for the job can make a huge difference.  
%Sometimes recognizing the best tool for the job can be tricky.  Of course, the more problems of
%this type you solve, the easier it gets.  Similarly, having more tools in your bag gives
%you more options.
%
%This also demonstrates something that is true of a lot of combinatorial problems:  There are often several valid ways of approaching them.  But there are also a lot of invalid approaches, so be careful!
%\end{rem}
%
%\begin{exer}
%In how many ways can $100$ be written as the sum of four positive
%integer summands?\\
%\answerlinetwo
%\begin{answer} We want the number of positive integer solutions to
%$$a + b + c + d = 100,$$ which by Theorem \ref{thm:distributions}
%is $$\binom{100-1}{4-1} =\binom{99}{3} = 156849.$$
%\end{answer}
%\end{exer}
%
%\noindent
%The following corollary is similar to Theorem \ref{thm:distributions} except that the numbers
%are allowed to be $0$.
%
%\begin{cor}\label{cor:demoivre}
%Let $n$ be a positive integer. The number of non-negative integer
%solutions to
%$$y_1 + y_2 + \cdots + y_k = n$$
%is
%$$\binom{n + k - 1}{k - 1}.$$
%\begin{pf}
%Set $x_i - 1 = y_i$ for $i=1,\ldots, k$. Then $x_i \geq 1$, and equation
%
%$$y_1 + y_2 + \cdots + y_k = n$$
%is equivalent to 
%$$x_1 - 1 + x_2 - 1 + \cdots + x_k - 1 = n$$ 
%which is equivalent to
%$$x_1 + x_2 + \cdots + x_k = n + k,$$
%which has $\binom{n + k - 1}{k - 1}$
%solutions by Theorem~\ref{thm:distributions}.
%\end{pf}
%\end{cor}
%
%Notice that the previous corollary is very similar to Theorem~\ref{thm:ballsbins}.
%We leave it for the interested reader to determine whether or not there is a
%deeper connection between these two.
%The next example demonstrates that the technique used in the proof of 
%Corollary~\ref{cor:demoivre} can be extended.
%
%\begin{exa}Find the number of quadruples $(a, b, c, d)$ of
%integers satisfying
%$$a + b + c + d  = 100, \ a \geq 30, b > 21, c \geq 1, d \geq 1.$$
%\begin{sol} Put $a' + 29 = a, b' +  20 = b.$ Then we want the number
%of positive integer solutions to
%$$a' + 29 + b' + 21 + c + d = 100,$$ or
%$$a' + b' + c + d = 50.$$
%By Theorem \ref{thm:distributions} this number is
%$$\binom{50-1}{4-1}=\binom{49}{3} = 18424.$$
%\end{sol}
%\end{exa}
%
%
%\begin{exer}
%In how many ways may $1024$ be written as the {\em product} of three
%positive integers? (Hint: Find the prime factorization of $1024$ and
%then figure out why that helps.)\\
%\answerlinethree
%\begin{answer} 
%Observe that $1024=2^{10}$. We need a decomposition of the
%form $2^{10} = 2^a 2^b2^c$, that is, we need integers solutions to
%$$a+b+c=10, \quad a\geq 0, b\geq 0, c \geq 0.$$
%By Corollary
%\ref{cor:demoivre} there are $\binom{10+3-1}{3-1} = \binom{12}{2} =
%66$ such solutions.
%\end{answer}
%\end{exer}

\newpage

\section{Binomial Theorem}\label{chap:bino_mio}
It is well known that
\begin{equation}(a + b)^2 = a^2 + 2ab + b^2\end{equation}Multiplying this last equality by
$a + b$ one obtains
$$(a + b)^3 = (a + b)^2(a + b) = a^3 + 3a^2b + 3ab^2 + b^3$$ Again, multiplying
\begin{equation}
(a + b)^3 = a^3 + 3a^2b + 3ab^2 + b^3
\end{equation}by $a + b$ one obtains
$$(a + b)^4 = (a + b)^3(a + b) = a^4 + 4a^3b + 6a^2b^2 + 4ab^3 + b^4$$

This generalizes, as we see in the next theorem.
\begin{thm}[Binomial Theorem]\index{Binomial Theorem}
Let $x$ and $y$ be variables and $n$ be a nonnegative integer.  Then
$$(x+y)^n = \sum_{i=0}^{n}\binom{n}{i}x^{n-i}y^i.$$
\end{thm}

\begin{exa}
Expand $(4x + 5)^3$, simplifying as much as possible.
\begin{sol}
\begin{eqnarray*}
(4x + 5)^3 & = & \binom{3}{0}(4x)^35^0 + \binom{3}{1}(4x)^2(5)^1 + \binom{3}{2}(4x)^1(5)^2 + \binom{3}{3}(4x)^05^3 \\
& = & (4x)^3 + 3(4x)^2(5) + 3(4x)(5)^2 + 5^3 \\
& = & 64x^3 + 240x^2 + 300x + 125
\end{eqnarray*}
\end{sol}
\end{exa}

\begin{exa}
In the following, $i=\sqrt{-1}$, so that $i^2=-1$.
$$
\begin{array}{lll}
(2 + i)^5 & = & 2^5 + 5(2)^4(i) + 10(2)^3(i)^2  + 10(2)^2(i)^3 + 5(2)(i)^4 + i^5 \\
& = & 32 + 80i - 80 - 40i + 10 + i \\
& = & -38 + 39i
\end{array}
$$
Notice that we skipped the step of explicitly writing out the binomial coefficient 
for this example.  You can do it either way--just make sure you aren't forgetting anything
or making algebra mistakes by taking shortcuts.
\end{exa}

\begin{exer}
Expand $(2x - y^2)^4$, simplifying as much as possible.
\vspace*{3in}
\begin{answer}
\begin{eqnarray*}
%\! is a small negative space. This is too wide, but barely and that fixes it.
(2x - y^2)^4 &\!\! = & \!\!\!\! \binom{4}{0}(2x)^4 + \binom{4}{1}(2x)^3(-y^2) + \binom{4}{2}(2x)^2(-y^2)^2 + \binom{4}{3}(2x)(-y^2)^3 + \binom{4}{4}(-y^2)^4 \\
 & = & (2x)^4 + 4(2x)^3(-y^2) + 6(2x)^2(-y^2)^2 + 4(2x)(-y^2)^3 + (-y^2)^4 \\
& = & 16x^4 - 32x^3y^2 + 24x^2y^4 - 8xy^6 + y^8
\end{eqnarray*}
\end{answer}
\end{exer}

The most important things to remember when using the binomial theorem are not to forget the binomial
coefficients, and not to forget that the powers (i.e. $x^{n-i}$ and $y^i$) apply to the whole term,
including any coefficients.  
A specific case that is easy to forget is a negative sign on the coefficient.
Did you make any of these mistakes when doing the last exercise?  Be sure to identify your
errors so you can avoid them in the future.

\begin{exer}
Expand $(\sqrt{3} + \sqrt{5})^4$, simplifying as much as possible.
\vspace*{3in}
\begin{answer}
$$
\begin{array}{lll}
(\sqrt{3} + \sqrt{5})^4 & = & (\sqrt{3})^4 +
4(\sqrt{3})^3(\sqrt{5})  + 6(\sqrt{3})^2(\sqrt{5})^2
+ 4(\sqrt{3})(\sqrt{5})^3 + (\sqrt{5})^4 \\
& = & 9 + 12\sqrt{15} + 90 +  20\sqrt{15} + 25 \\
& = & 124 + 32\sqrt{15}
\end{array}
$$
\end{answer}
\end{exer}

\begin{exa}Let $n \geq 1$.
Find a closed form for $\dsum_{k=0}^{n} \binom{n}{k}(-1)^k$.
\begin{sol}
Using a little algebra and the binomial theorem,  we can see that 
$$\dsum_{k=0}^{n} \binom{n}{k}(-1)^k
=\dsum_{k=0}^{n} \binom{n}{k}1^{n-k}(-1)^k
=(1-1)^n
=0.$$

\vspace*{0in}
\end{sol}
\end{exa}

\begin{exer}
Find a closed form for $\dsum_{k=1}^{n} \binom{n}{k}3^k$.
\vspace*{2in}
\begin{answer}
Using a little algebra and the binomial theorem,  we can see that 
$$\dsum_{k=1}^{n} \binom{n}{k}3^k
=\dsum_{k=0}^{n} \binom{n}{k}3^{k}-1
=\dsum_{k=0}^{n} \binom{n}{k}1^{n-k}3^{k}-1
=(1+3)^{n}-1=4^n-1.$$
\end{answer}
\end{exer}

If we ignore the variables in the Binomial Theorem and write down the coefficients
for increasing values of $n$, a pattern, called {\em Pascal's Triangle}, emerges
(see Figure~\ref{fig:pascal}).
\index{Pascal's Triangle}

\begin{figure}[h]
\centering
\small{1 \\
1 \hspace{3mm} 1\\
1 \hspace{3mm} 2 \hspace{3mm} 1 \\
1 \hspace{3mm} 3 \hspace{3mm} 3 \hspace{3mm} 1 \\
1 \hspace{3mm} 4 \hspace{3mm} 6 \hspace{3mm} 4 \hspace{3mm} 1 \\
1 \hspace{3mm} 5 \hspace{3mm} 10 \hspace{3mm} 10 \hspace{3mm} 5 \hspace{3mm} 1 \\
1 \hspace{3mm} 6 \hspace{3mm} 15 \hspace{3mm} 20 \hspace{3mm} 15 \hspace{3mm} 6 \hspace{3mm} 1 \\
1 \hspace{3mm} 7 \hspace{3mm} 21 \hspace{3mm} 35 \hspace{3mm} 35 \hspace{3mm} 21 \hspace{3mm} 7 \hspace{3mm} 1 \\
1 \hspace{3mm} 8 \hspace{3mm} 28 \hspace{3mm} 56 \hspace{3mm} 70 \hspace{3mm} 56 \hspace{3mm} 28 \hspace{3mm} 8 \hspace{3mm} 1\\
1 \hspace{3mm} 9 \hspace{3mm} 36 \hspace{3mm} 84 \hspace{3mm} 126 \hspace{3mm} 126 \hspace{3mm} 84 \hspace{3mm} 36 \hspace{3mm} 9 \hspace{3mm} 1 \\
1 \hspace{3mm} 10 \hspace{3mm} 45 \hspace{3mm} 120 \hspace{3mm} 210 \hspace{3mm} 252 \hspace{3mm} 210 \hspace{3mm} 120 \hspace{3mm} 45 \hspace{3mm} 10 \hspace{3mm} 1 \\
$\vdots$
}\caption{Pascal's Triangle}\label{fig:pascal}\index{Pascal's Triangle}
\end{figure}

Notice that each entry (except for the $1$s) is the sum of the two neighbors just above it.
This observation leads to the Pascal's Identity.

\begin{thm}[Pascal's Identity]\label{thm:pascalIdent}\index{Pascal's Identity}
Let $n$ and $k$ be positive integers with $k\leq n$.  Then
$$\binom{n+1}{k}=\binom{n}{k-1}+\binom{n}{k}.$$
\begin{pf}
See Problem~\ref{pro:pascident}.
\end{pf}
\end{thm}

%%%%%%%%%%%%%%%%%%%%%%%%%%%%%%%%%%%%%%%%

\section{Inclusion-Exclusion}\label{sec:inclExcl}
The Sum Rule (Theorem~\ref{rul:sum}) gives us the cardinality for unions of
finite sets that are mutually disjoint. In this section we will drop
the disjointness requirement and obtain a formula for the
cardinality of unions of general finite sets.

The Principle of {\em Inclusion-Exclusion} is attributed to both
Sylvester and to Poincar\'{e}.  We will only consider the cases involving
two and three sets, although the principle easily generalizes to $k$ sets.

\begin{thm}[Inclusion-Exclusion for Two Sets]\index{inclusion-exclusion!two sets}
\label{thm:two_set_incl_excl}
Let $A$ and $B$ be sets.  Then 
$$|A \cup B| = |A| + |B| - |A \cap B|$$

\begin{pf}
Clearly there are $|A\cap B|$ elements that are in both $A$ and $B$.
Therefore, $|A|+|B|$ is the number of element in $A$ and $B$, where
the elements in $|A\cap B|$ are counted twice.  From this it is clear
that $|A \cup B| = |A| + |B| - |A \cap B|$.
\end{pf}
\end{thm}
%
%\begin{figure}[h]
%\begin{minipage}{.45\textwidth}
%$$
%\psset{unit=1pc} \pscircle(-1,0){2} \pscircle(1,0){2}
%\rput(0,0){\tiny R_1} \rput(-2, 0){\tiny R_2} \rput(2, 0){\tiny R_3}
%\rput(-3, 1.5){\tiny A} \rput(3, 1.5){\tiny B} $$ \caption{Two-set Inclusion-Exclusion}
%\label{fig:2_set_incl_excl}
%\end{minipage}
%\begin{minipage}{.45\textwidth}
%$$
%\psset{unit=1pc} \pscircle(-1,0){2} \pscircle(1,0){2}
%\rput(0,0){\tiny 8} \rput(-2, 0){\tiny 18} \rput(2, 0){\tiny 6}
%\rput(2.5, 2.5){\tiny 6} \rput(-3, 1.5){\tiny A} \rput(3, 1.5){\tiny
%B}$$
%\caption{Example
%\ref{exa:incl_excl_1}.}\label{fig:incl_excl_1}\end{minipage}
%\end{figure}

\begin{exa}\label{exa:incl_excl_1}
Of  $40$ people, $28$ smoke and  $16$ chew tobacco. It is also known
that  $10$ both smoke and chew. How many among the $40$ neither
smoke nor chew?  
\begin{sol} Let $A$
denote the set of smokers and $B$ the set of chewers. Then
$$|A \cup B| = |A| + |B| - |A \cap B| = 28 + 16 - 10 = 34,$$
meaning that there are 34 people that either smoke or chew (or
possibly both). Therefore the number of people that neither smoke
nor chew is $40 - 34 = 6.$
\end{sol}
\end{exa}



\begin{exer}
In a group of $100$ camels, $46$ eat wheat, $57$ eat barley, and
$10$ eat neither. How many camels eat both wheat and barley?
\vspace*{3in}
\begin{answer} 
Let $A$ be the set of camels eating wheat and $B$ be the set of camels eating barley.
We know that $|A|=46$, $|B|=57$, and $|A\cup B|=100-10=90$.
We want $|A\cap B|$.  By Theorem~\ref{thm:two_set_incl_excl} (solving it for $|A\cap B|$),
$$|A\cap B| = |A|+|B|-|A \cup B| = 46+57-90=13.$$
\end{answer}
\end{exer}

\begin{exa}
Consider the set $A$ that are multiples of $2$ no greater than 114.  That is, 
$$A = \{2,4,6, \ldots , 114\}.$$
\begin{lenum}
\item 
How many elements are there in $A$?\item How many are
divisible by $3$? \item How many are divisible by $5$?  \item How
many are divisible by $15$? \item How many are divisible by either
$3$, $5$ or both? \item How many are neither divisible by $3$ nor
$5$? 
\item 
How many are divisible by exactly one of $3$ or $5$?
\end{lenum}
\begin{sol} Let $A_k \subset A$ be the set of those integers divisible by $k$.
\begin{lenum}
\item 
Notice that the elements are $2 = 2(1)$, $4 = 2(2)$, \ldots
, $114 = 2(57)$. Thus $|A| = 57.$ 
\item 
Notice that 
$$A_3=\{6, 12, 18, \ldots , 114\}=\{1\cdot 6,2\cdot 6, 3\cdot 6,\ldots, 19\cdot 6 \},$$ 
so $|A_3| = 19$.  
\item 
Notice that 
$$A_5=\{10, 20, 30,\ldots , 110\}=\{1\cdot 10,2\cdot 10, 3\cdot 10,\ldots, 11\cdot 10 \},$$ 
so $|A_5| = 11$.  
\item Notice that  
$A_{15}=\{30, 60, 90\}$, so $|A_{15}|=3$.
\item 
First notice that $A_3\cap A_5=A_{15}$.  Then it is clear that
the answer is 
$$|A_3 \cup A_5| = |A_3|+|A_5|-|A_{15}| = 19 + 11 -3= 27.$$ 
\item
We want
$$|A \setminus (A_3 \cup A_5)| = |A| - |A_3 \cup A_5| = 57 - 27 = 30.$$ 
\item We want
$$\begin{array}{lll}|(A_3 \cup A_5) \setminus (A_3 \cap A_5)|
& =  & |(A_3 \cup A_5)| - |A_3 \cap A_5|\\ &  = &  27 - 3
\\ & = &  24.\end{array}$$
\end{lenum}
\end{sol}
\end{exa}

%We will use the following somewhat intuitive result in several examples.
%\begin{thm}
%Let $k$ and $n$ be a positive integers with $k<n$.  Then there are
%$\left\lfloor\frac{n}{k}\right\rfloor$ numbers between $1$ and $n$, inclusive,
%that are divisible by $k$.
%\end{thm}
%
%\begin{exer}
%How many integers between $1$ and $1000$ inclusive, do not share a
%common factor with $1000$, that is, are relatively prime to $1000$?
%(Hint: $1000$ only has 2 prime factors.  
%Start by using inclusion/exclusion
%to count the number of numbers that have either of these as a factor.)\\[3.75in]
%\begin{answer}   
%Observe that $1000 = 2^35^3$, and thus from the $1000$
%integers we must weed out those that have a factor of $2$ or of $5$
%in their prime factorization. If $A_2$ denotes the set of those
%integers divisible by 2 in the interval $[1, 1000]$ then clearly
%$\dis{ |A_2| = \left\lfloor \frac{1000}{2} \right\rfloor} = 500$.
%Similarly, if $A_5$ denotes the set of those integers divisible by
%$5$ then $\dis{ |A_5| = \left\lfloor \frac{1000}{5} \right\rfloor} = 200$.
%Also $\dis{ |A_2 \cap A_5| = \left\lfloor \frac{1000}{10} \right\rfloor} =
%100$. This means that there are $|A_2 \cup A_5| = 500 + 200 -
%100 = 600$ integers  in the interval $[1, 1000]$ sharing at least a
%factor with $1000$, thus there are $1000 - 600 = 400$ integers in
%$[1, 1000]$ that do not share a factor prime factor with $1000$.
%\end{answer}
%\end{exer}

We now derive a three-set version of the Principle of
Inclusion-Exclusion.

\newpage

\begin{thm}[Inclusion-Exclusion for Three Sets]\label{thm:three_set_incl_excl}\index{inclusion-exclusion!three sets}
Let $A$, $B$, and $C$ be sets.  Then 
$$
\begin{array}{lll}
|A \cup B \cup C| & = & |A| + |B| + |C|\\ & & -
|A \cap B| - |A \cap C |   - |B \cap C| \\ & & +
|A \cap B \cap C|
\end{array}
$$

\noindent
{\bf Proof:}
Using the associativity and distributivity of unions of sets, we see
that
\begin{eqnarray*}
|A \cup B \cup C| & = & |A \cup (B \cup C)| \\
& = & |A| + |B \cup C| - |A \cap (B \cup C)| \\
& = & |A| + |B \cup C| - |(A \cap B) \cup (A \cap C)|  \\
& = & |A| + |B| + |C| - |B \cap C| - |A \cap B| - |A \cap C|+ |(A\cap B)\cap (A \cap C)| \\
& = & |A| + |B| + |C| - |B \cap C| - \left( |A \cap B| + |A \cap C|  - |A\cap B \cap C|\right) \\
& = & |A| + |B| + |C|  - |A \cap B|- |B \cap C | - |C \cap A|  + |A \cap B\cap C|. \, \, \square
\end{eqnarray*}
\end{thm}


\begin{exa} 
At {\em Medieval High} there are forty students. Amongst
them, fourteen like Mathematics, sixteen like theology, and eleven
like alchemy. It is also known that seven like Mathematics and
theology, eight like theology and alchemy and five like Mathematics
and alchemy. All three subjects are favored by four students. How
many students like neither Mathematics, nor theology, nor alchemy?

\bigskip
\noindent
{\bf Solution:}   
Let $A$ be the set of students liking Mathematics, $B$ the set
of students liking theology, and $C$ be the set of students liking
alchemy. We are given that $$ |A| = 14, |B| = 16, |C|
= 11, |A\cap B| = 7, |B\cap C| = 8, |A\cap C| = 5,$$
and
$$ |A\cap B\cap C| = 4.$$
Using Theorem~\ref{thm:three_set_incl_excl}, along with some set identities, 
we can see that
\begin{eqnarray*}
|\overline{A} \cap \overline{B} \cap \overline{C}| 
&=& |\overline{A\cup B\cup C}|\\
&=& |U| - |A\cup B\cup C|\\
&=&|U| - |A| - |B| - |C| + |A\cap B| + |A\cap C| +
|B\cap C| - |A\cap B \cap C|\\
&=& 40 - 14 - 16 - 11 + 7 + 5 + 8 - 4\\
&=&15.\
\end{eqnarray*}
\end{exa}

\newpage

\begin{exer}
A survey of a group's viewing habits revealed 
the percentages that watch a given sport.  The results are given below.
Calculate the percentage of the group that watched none of the three sports.

\begin{tabular}{p{.25in}lll}
&$28\%$  gymnastics & $14\%$  gymnastics \& baseball&  $8\%$  all three sports\\
&$29\%$  baseball& 		$10\%$  gymnastics \& soccer\\
&$19\%$  soccer & 		$12\%$  baseball \& soccer 
\end{tabular}

\vspace*{3.25in}
\begin{answer}
Using Theorem~\ref{thm:three_set_incl_excl}, we know that
$28+29+19-14-10-12+8=48\%$  watch at least one of these sports.  
That leaves $52\%$ that don't watch any of them.
\end{answer}
\end{exer}

%
%\begin{exa}
%How many integers between $1$ and $600$ inclusive are not divisible
%by $3,$ nor $5$, nor $7$?
%\begin{sol} Let $A_k$ denote the numbers in $[1, 600]$ which are
%divisible by $k = 3, 5, 7.$ Then
%
%%\renewcommand{\arraystretch}{2}
%\begin{eqnarray*}
%|A_3| & = & \left\lfloor\frac{600}{3}\right\rfloor         =  200, \\
%|A_5| & = & \left\lfloor\frac{600}{5}\right\rfloor         =  120, \\
%|A_7| & = & \left\lfloor \frac{600}{7}\right\rfloor        =  85, \\
%|A_{15}| &  = & \left\lfloor \frac{600}{15}\right\rfloor   =  40, \\
%|A_{21}| &  = & \left\lfloor \frac{600}{21}\right\rfloor   =  28,\\
%|A_{35}|&  = & \left\lfloor \frac{600}{35}\right\rfloor   =  17, \mbox{ and} \\
%|A_{105}|&  = &\left\lfloor \frac{600}{105}\right\rfloor  =  5.
%\end{eqnarray*}
%By Inclusion-Exclusion there are $200 + 120 + 85 - 40 - 28 - 17 + 5
%= 325$ integers in $[1, 600]$ divisible by at least one of 3, 5, or
%7. Those not divisible by these numbers are a total of $600 - 325 =
%275.$
%\end{sol}
%\end{exa}


\begin{exer}
Would you believe a market investigator that reports that of $1000$
people, $816$ like candy, $723$ like ice cream, $645$ like cake, while
$562$ like both candy and ice cream, $463$ like both candy and cake,
$470$ like both ice cream and cake, while $310$ like all three? State
your reasons! 
\vspace*{3in}
\begin{answer} 
Let $C$ denote the set of people
who like candy, $I$ the set of people who like ice cream, and $K$
denote the set of people who like cake. We are given that $|C|
= 816$, $|I| = 723$, $|K| = 645$, $|C\cap I| = 562$,
$|C\cap K| = 463$, $|I\cap K| = 470$, and $|C\cap I\cap K| = 310$. By Inclusion-Exclusion we have
$$\begin{array}{lll}
|C \cup I \cup K| & = & |C| + |I| + |K| \\ & & -
|C \cap I| - |C \cap K |
- |I \cap C| \\ & &  + |C \cap I \cap K| \\
& = & 816 + 723 + 645 - 562 - 463 - 470 + 310 \\ & = & 999.
\end{array}$$
The investigator miscounted, or probably did not report one person
who may not have liked any of the three things.
\end{answer}
\end{exer}

% Not that interesting for CS students.
% CSC, 6/18/14
%\begin{exa}\item How many integral solutions to the equation
%$$a + b + c + d = 100,$$are there given the following
%constraints:
%$$1 \leq a \leq 10,\ b \geq 0, \ c \geq 2, 20 \leq d \leq 30 ?$$
%\begin{sol} We use Inclusion-Exclusion.
%There are $\binom{80}{3} =82160$ integral solutions to $$a + b + c +
%d = 100, \ \ a \geq 1, b \geq 0, c \geq 2, d \geq 20.$$ Let $A$ be
%the set of solutions with $$ a \geq 11, b \geq 0, c \geq 2, d \geq
%20$$ and $B$ be the set of solutions with
%$$ a \geq 1, b \geq 0, c \geq 2, d \geq
%31.$$ Then $|A| = \binom{70}{3}$, $|B| = \binom{69}{3}$,
%$|A \cap B| =  \binom{59}{3}$ and so
%$$|A \cup B| = \binom{70}{3} + \binom{69}{3} - \binom{59}{3} = 74625.$$
%The total number of solutions to
%$$a + b + c + d = 100$$with
%$$1 \leq a \leq 10,\ b \geq 0, \ c \geq 2, 20 \leq d \leq 30 $$
%is thus
%$$\binom{80}{3} - \binom{70}{3} - \binom{69}{3} + \binom{59}{3} = 7535.$$
%\end{sol}
%\end{exa}

\begin{exa}
An auto insurance company has $10,000$ policyholders. Each policy
holder is classified as
\begin{center} 
\begin{itemize}
\item young or old, 
\item male or female, and 
\item married or single. 
\end{itemize}
\end{center}
Of these policyholders, 3000 are young, 4600 are male, and 7000 are
married. The policyholders can also be classified as 1320 young
males, 3010 married males, and 1400 young married persons. Finally,
600 of the policyholders are young married males.

How many of the company's policyholders are young, female, and
single? 
\begin{sol}
Let $Y, F, S, M$ stand for young, female,
single, male, respectively, and let $Ma$ stand for married. We have
$$\begin{array}{lll}|Y\cap F \cap S| & = & |Y\cap F| - |Y\cap F \cap Ma| \\
& =  &|Y| - |Y\cap M|  \\
& & \qquad - (|Y\cap Ma|  - |Y\cap Ma \cap M|) \\
& = & 3000 - 1320 - (1400 - 600) \\
& = & 880.
\end{array}
$$
\end{sol}
\end{exa}

The following problem is a little more challenging than the others we have seen, but
you have all of the tools you need to tackle it.

\begin{exer}[Lewis Carroll in {\em A Tangled Tale.}]
In a very hotly fought battle, at least $70\%$ of the combatants
lost an eye, at least $75\%$ an ear, at least $80\%$ an arm, and at
least $85\%$ a leg. What can be said about the percentage who lost
all four members? 
\vspace*{3.5in}
\begin{answer} 
We can either use inclusion-exclusion for four sets or use a few applications
of inclusion-exclusion for two sets.  Let's try the latter.

Let $A$ denote the set of
those who lost an eye, $B$ denote those who lost an ear, $C$ denote
those who lost an arm and $D$ denote those losing a leg. Suppose
there are $n$ combatants.
 Then
$$\begin{array}{lll}
n & \geq & | A \cup B| \\ & = &  | A| + | B| - | A \cap B| \\ & = & .7n + .75n - | A\cap B|, 
\end{array}$$
$$\begin{array}{lll}
n  &\geq & | C \cup D| \\ & = & | C| + | D| - | C \cap D| \\ & = & .8n + .85n - |C\cap D|. 
\end{array}$$
This gives
$$| A \cap B| \geq .45n,$$
$$| C \cap D| \geq .65n.$$This means that
$$\begin{array}{lll}n &\geq & | (A\cap B) \cup (C \cap D)| \\ & = & | A \cap B| + | C \cap D| - | A \cap B \cap C \cap D|\\
& \geq &  .45n + .65n - | A \cap B \cap C \cap D|,\end{array}$$
whence
$$| A \cap B \cap C \cap D| \geq .45 + .65n - n = .1n.$$This means
that at least $10\%$ of the combatants lost all four members.
\end{answer}
\end{exer}

%\newpage